\section{Introduction}
\label{sec:intro}

When you write a passage, a large language model that reads it can often tell \emph{how} you produced it. It can usually tell whether the text was dictated rather than typed, whether another LLM wrote it, and --- in principle --- which keyboard layout shaped the pattern of your typos. We call each such latent variable a \modes of text production. The question this paper asks is whether different production modes share a common geometry inside an LLM, or whether each lives in its own corner of the representation space.

{\bf Why does this matter?} The same axis matters to several research communities that have so far worked separately on it. AI-text detection treats the AI-vs-human axis as a standalone binary classification problem \citep{mitchell2023detectgpt,verma2023ghostbuster,schafer2025stylistic}. Authorship attribution and stylometry treat \emph{who wrote it} as another standalone classification problem with its own decades-old methodology \citep{rivera2021universal,sarfati2025prompt}. Interpretability work on linear probes and sparse autoencoders \citep{alain2016linear,templeton2024scaling} shows that high-level features can be linearly decoded from hidden states, but tends to study one feature at a time. Nobody has asked: do these various \modes coexist in the same residual stream, and are they geometrically disentangled?

{\bf The gap.} A unified picture would carry practical weight. If keyboard layout, dictation status, or AI-vs-human origin are all recoverable from the same hidden state, then any system that runs an LLM over user text leaks more than the text. If \modes are orthogonal, activation steering along one mode (\eg ``make this less formal'') should not disturb another (``preserve the author's voice''). And if AI text occupies a separate latent direction from the actual content of a passage, benchmark contamination by AI generations is detectable in hidden states even when invisible at the surface. None of these claims is testable without first establishing whether \modes are jointly represented and whether they share or separate directions.

{\bf Our approach.} We treat \modes as a unified latent variable and probe four very different ones --- 7-way literary authorship, AI-vs-human authorship, formal-vs-casual register, and QWERTY-vs-Dvorak typo patterns --- in the last-token hidden states of \qwen. For each \modes, at each layer $L \in \{0, 2, \ldots, 28\}$ and each context length $N \in \{32, 64, 128\}$, we fit a linear probe under 5-fold stratified cross-validation. We then ask three geometric questions: (i) how do the probe directions for different \modes align with each other (cosine similarity)? (ii) does a probe direction trained on one \modes carry information about another (cross-mode transfer)? (iii) how does each \modes emerge as a function of network depth?

{\bf Preview of results.} A single linear probe decodes every \modes well above chance and well above random-direction baselines, with peak accuracies of \textbf{79.2\%--99.8\%}. The decoded directions are pairwise near-orthogonal --- all six pairwise cosines satisfy $|\cos| \leq 0.18$. Cross-mode transfer is near chance for orthogonal modes (\typo $\leftrightarrow$ \register at 0.48--0.50) but reveals one expected content overlap: the formality direction recovers 74\% of AI-vs-human labels, the latent footprint of ChatGPT's well-documented formality bias. Layer-depth profiles further split the four \modes into three tiers: content modes (register, AI-vs-human) saturate by layer 2, authorship rises gradually across the entire stack, and surface-typing artifacts become decodable only at the final layers.

{\bf Contributions.} We make the following contributions:
\begin{itemize}[leftmargin=*,itemsep=0pt,topsep=0pt]
    \item We propose the \modes abstraction --- production mode as a unified latent variable --- and instantiate it on four diverse axes spanning content, author identity, register, and surface-typing artifact.
    \item We show that a single linear probe family on \qwen last-token hidden states decodes each \modes at 79\%--99\% accuracy, establishing decodability as a property shared across very different mode types.
    \item We characterise the \emph{geometry} of \modes: pairwise direction cosines $\leq 0.18$ and cross-mode probe transfer near chance for independent modes, providing the first quantitative evidence that production modes are largely disentangled in LLM hidden states.
    \item We identify a three-tier layer-depth taxonomy --- shallow content modes, gradual author modes, and late surface modes --- that organises a previously unstructured set of probe findings.
\end{itemize}

\Secref{sec:related} surveys prior work on AI-text detection, stylometry, and probing. \Secref{sec:method} describes our model, corpora, probes, and geometric analyses. \Secref{sec:results} presents quantitative results. \Secref{sec:discussion} interprets the geometry and notes limitations. \Secref{sec:conclusion} summarises and lists follow-up directions.
